\section{Discussion}

Feature normalization changes how quickly a model learns different coefficients
under a finite training budget. We have shown that this can affect the best
attainable rare-subgroup risk even when each method may choose any nonnegative
constant learning rate and any deterministic stop within the budget.
Standardization has lower risk on the specified rare-signal family, whereas
unscaled training has lower risk on the common-signal family. The gaps remain
bounded away from zero as prevalence decreases while the expected number of
rare observations stays of order one. The theorem thus gives sufficient
conditions for both orderings, although it does not locate a general boundary
between them.

The supporting comparisons explain how standardization can reduce bias in the
rare coefficient while increasing noise, and how dependence can make aggregate
and rare-presence prediction favor different rules. We describe both effects
through the same error moments. To interpret these comparisons, we need to
specify the learning rate, batch size, stopping count, and prediction objective
together with the normalization rule.

The experiments show that an advantage under a fixed calibration can disappear
after tuning. Both original anchor comparisons disappear under the declared
oracle searches. We still observe differences in relative learning with
estimated preprocessing and validation selection, but absent features and noisy
validation add uncertainty. On Mushroom, expanding the rate search removes the
resolved large-sample rare-risk advantage of SD scaling under aggregate
selection, while damping retains its advantage. The four-rule extension removes
the remaining rate boundaries on this dataset. Variance scaling has higher rare
risk in the primary comparison, while its aggregate difference remains
unresolved. Some synthetic selections still reach rate boundaries, and numerical
checks disagree on extreme paths.

The separate damping comparison finds improvements at several fixed offsets,
but also a severe failure under rare validation in the smaller sample. The
original offset $0.01$ remains a heuristic. Choosing the offset would require
additional tuning effort, which should be included in the comparison. More
generally, validation cannot directly assess rare-group performance when the
group is absent from the validation data.

Our theoretical comparison assumes two independent features, known population
preprocessing, a fixed intercept, one fresh observation per update, and the
stated signal, prevalence, and budget ranges. It covers all nonnegative constant
rates, while the empirical studies use finite searches. Extending the theory to
estimated scales and validation selection remains open. The results concern
fixed diagonal scaling, and the evolving preconditioners of adaptive optimizers
require a separate analysis. The next empirical step is to declare the predicted
orderings before evaluating new synthetic
configurations and an additional categorical task. To support a broader
normalization recommendation, we would also need to analyze how selection
depends on rare-group coverage.
